\section{Contexto y repaso del curso}

Esta clase es la novena del curso KAIST CS492D (Modelos de difusión y modelos generativos), con el tema \textbf{Flow Matching} (flujo coincidente). En las clases anteriores, discutimos en detalle las diversas variantes de los modelos de difusión —desde DDPM hasta DDIM, pasando por modelos basados en SDE/ODE y técnicas de aceleración como DPM-Solver—. Esta clase introducirá los modelos generativos desde una perspectiva totalmente nueva: \textbf{Continuous Normalizing Flows} (flujos normalizadores continuos) y \textbf{Flow Matching} (flujo coincidente).

\begin{knowledgebox}{Repaso de la estructura del curso}
En las clases anteriores, establecimos la siguiente cadena de conocimientos:
\begin{enumerate}
    \item \textbf{DDPM}: modelo de difusión de desruido en pasos discretos, que típicamente requiere 1000 pasos de muestreo.
    \item \textbf{DDIM}: muestreo determinístico, que reduce los pasos a 50--200.
    \item \textbf{Score-Based SDE}: unifica el proceso de difusión bajo un marco de ecuaciones diferenciales estocásticas (SDE).
    \item \textbf{Probability Flow ODE}: deriva una forma equivalente en ODE desde la SDE, logrando generación determinística.
    \item \textbf{DPM-Solver}: un solver especializado que aprovecha la estructura de ODE, reduciendo aún más los pasos a 10--15.
\end{enumerate}
El Flow Matching de esta clase ofrecerá otra forma de comprender y entrenar modelos generativos basados en ODE.
\end{knowledgebox}

\subsection{Cuello de botella de velocidad de los modelos de difusión}

La ventaja principal de los modelos de difusión radica en su capacidad para generar salidas de alta calidad sin colapso de modos (mode collapse). Sin embargo, su cuello de botella principal sigue siendo la \textbf{velocidad de muestreo}:

\begin{table}[H]
\centering
\begin{tabular}{lcc}
\toprule
\textbf{Método} & \textbf{Pasos típicos} & \textbf{¿Requiere reentrenamiento?} \\
\midrule
DDPM & $\sim$1000 & No \\
DDIM & 50--200 & No \\
DPM-Solver (orden superior) & 10--15 & No \\
Consistency Models & 1--2 & Sí (fine-tuning) \\
Métodos de destilación & 1--5 & Sí (fine-tuning) \\
\bottomrule
\end{tabular}
\caption{Comparativa de métodos de aceleración del muestreo en modelos de difusión}
\end{table}

El ponente señala que la evolución desde DDPM hasta DPM-Solver logró una aceleración de aproximadamente \textbf{100 veces}, sin necesidad de reentrenar el modelo —una característica muy atractiva de los modelos de difusión—. Sin embargo, para reducir aún más los pasos a 1--5, se requieren métodos adicionales de entrenamiento como la destilación.

\subsection{Introducción a Consistency Models}

\begin{importantbox}{Idea central de Consistency Models}
La idea básica de Consistency Models es: entrenar una red para que pueda predecir la salida final directamente desde \textbf{cualquier punto intermedio} en la trayectoria de ODE. Específicamente:
\begin{enumerate}
    \item Partiendo de una muestra gaussiana estándar, se ejecuta un solver de ODE para obtener la trayectoria completa.
    \item Para cada punto intermedio $x_t$ en la trayectoria, se calcula la expectativa de la salida final utilizando la fórmula de Tweedie.
    \item Se entrena la red para que las predicciones en momentos tempranos sean coherentes con las predicciones en momentos tardíos (consistent).
\end{enumerate}
De este modo, durante la inferencia, basta con uno o dos pasos para saltar directamente desde el ruido hasta la muestra de datos.
\end{importantbox}

El inconveniente común tanto de Consistency Models como de otros métodos de destilación es que ambos requieren un modelo de difusión preentrenado como punto de partida, logrando la aceleración mediante un ajuste adicional (fine-tuning). Esto implica dedicar más recursos computacionales durante la fase de entrenamiento a cambio de una velocidad mayor en tiempo de prueba.

\subsection{Resumen de la sección}

Desde DDPM hasta Consistency Models, la comunidad de modelos de difusión ha logrado avances significativos en la aceleración del muestreo. No obstante, estos métodos se ven limitados ya sea por un cronograma de ruido (noise schedule) específico o requieren entrenamiento de destilación complejo. Flow Matching abordará el problema desde una perspectiva completamente diferente, ofreciendo un marco más flexible y conciso para entrenar modelos generativos.

